% Source: https://artofproblemsolving.com/community/c3h1904914_area_of_fourth_region
% Stable identity: aops-topic-1904914-q01
(a) Join $P$ also to the four vertices; this splits each quadrilateral into two triangles. On side $AB$ with midpoint $M_1$, triangles $PAM_1$ and $PM_1B$ have equal bases $AM_1=M_1B$ and the same height (the distance from $P$ to $AB$), so they have equal areas; call this area $p$. In the same way the two triangles on side $BC$ each have area $q$, those on $CD$ area $r$, and those on $DA$ area $s$.

\begin{asy}
usepackage("amsmath");
size(6.5cm);
defaultpen(fontsize(10pt));
pair A=(0,0), B=(8,0), C=(8,8), D=(0,8), P=(5,6);
pair M1=(4,0), M2=(8,4), M3=(4,8), M4=(0,4);
draw(A--B--C--D--cycle,linewidth(0.9));
draw(P--M1);draw(P--M2);draw(P--M3);draw(P--M4);
dot(P);dot(M1);dot(M2);dot(M3);dot(M4);
label("$A$",A,SW);label("$B$",B,SE);label("$C$",C,NE);label("$D$",D,NW);
label("$P$",P,E);
draw(P--A,dashed);draw(P--B,dashed);draw(P--C,dashed);draw(P--D,dashed);
pen q=blue+fontsize(9pt);
label("$p$",(3,2),q);label("$p$",(5.67,2),q);
label("$q$",(7,3.33),q);label("$q$",(7,6),q);
label("$r$",(5.67,7.33),q);label("$r$",(3,7.33),q);
label("$s$",(1.67,6),q);label("$s$",(1.67,3.33),q);
label("\mbox{}",(min(currentpicture,user=true)+max(currentpicture,user=true))/2,black); // keeps the text after the figure black when LaTeX typesets the labels inline
\end{asy}

The quadrilateral at each vertex consists of one triangle from each of the two sides meeting there:
\[
[A]=p+s,\qquad [B]=p+q,\qquad [C]=q+r,\qquad [D]=r+s .
\]
Hence $[A]+[C]=p+q+r+s=[B]+[D]$; each sum is half the area of the square. $\blacksquare$

(b) By (a), $[D]=[A]+[C]-[B]=22+10-18=\boxed{14}$. Such a configuration exists: for a square of side $8$ with $A=(0,0)$ and $P=(5,6)$ we get $p=12$, $q=6$, $r=4$, $s=10$ and the areas $22,18,10,14$.
